\section{Introduction}
\label{sec:intro}

% Teaser figure: generated by paper/submission_folder/figures/src/teaser.py from completed runs.
\begin{figure}[t]
  \centering
  \IfFileExists{submission_folder/figures/teaser.pdf}{%
    \includegraphics[width=\linewidth]{teaser.pdf}}{%
    \fbox{\parbox[c][4.2cm][c]{0.96\linewidth}{\centering\small\textcolor{mutedgray}{Teaser figure pending completed runs}}}}
  \caption{\textbf{\ours{} forecasts by moving observed features.} Given the last observed frame
  of real robot and surgical video and a sequence of future actions, \ours{} predicts, for every future
  step, where each patch's features come from (arrows: learned transport field), where new content must
  be synthesised (gate), and a correction. Recursive predictors regenerate the full feature grid at
  every step and drift (right).}
  \label{fig:teaser}
\end{figure}

A world model predicts how observations change in response to actions. Recent work shows that
forecasting in the feature space of a frozen, self-supervised encoder such as DINOv2 or V-JEPA~2
gives compact models that support planning without generating pixels
\citep{zhou2025dinowm,assran2025vjepa2,maes2026lewm,sobal2025pldm}. These predictors are
\emph{generative in feature space}. At every future step they output the complete grid of patch
features, usually by feeding their own predictions back as inputs. In manipulation and surgical
video, however, most of the scene persists: tissue, tools, the table and the robot move, but they
are rarely created or destroyed. Regenerating persistent content at every step wastes capacity and
lets errors compound under recursion.

We ask a simple question. \emph{Can a latent world model forecast by transporting features it has
already observed, and synthesise only what transport cannot explain?} Motion-based video
prediction answered a version of this question in pixel space
\citep{finn2016advection,ebert2017skip,debrabandere2016dfn,liu2017voxelflow}. The idea has not been
tested for foundation-feature world models used for planning, where features are semantic,
high-dimensional and coarse in space. Three questions are open: whether transport helps once the
predictor is strong, whether it survives held-out real data, and what it costs for control.

We introduce \ours{} (\cref{fig:teaser,fig:method}). A memory encoder reads the patch features of
the last few observed frames. For each future step, a decoder conditioned on the causal
action-prefix state produces per-patch queries. Each query (i) attends to a local window of
\emph{observed} feature locations and copies their content, (ii) gates between staying in place and
moving, and (iii) adds a learned correction. Every horizon is predicted directly from measured
features, so errors do not compound through recursion, and the transport weights form a motion
field that we can inspect and compare against optical flow.

\paragraph{Contributions.}
\begin{itemize}[leftmargin=1.2em,itemsep=1pt,topsep=2pt]
  \item \textbf{Method.} \ours{}, an action-conditioned local feature transport decoder for latent
  world models. It is a drop-in output head for any frozen patch encoder and adds $<\!2\%$ parameters.
  \item \textbf{Evidence across domains.} Matched comparisons against recursive (DINO-WM-style),
  direct-additive and cross-attention predictors on held-out real robot video (DROID), seven surgical
  dVRK tasks (Open-H-Embodiment), IWS manipulation forecasting and three planning benchmarks. We use
  two encoder families and report forecasting, decoded-pixel, probing, action-sensitivity and planning
  metrics.
  \item \textbf{Analysis.} We show that the learned transport field tracks measured motion. We also
  characterise a trade-off specific to anchored predictors: open-loop accuracy against action
  sensitivity. An action-contrastive objective resolves this trade-off for planning.
\end{itemize}
